\documentclass{article}
\usepackage{graphicx} % Required for inserting images
\usepackage{amsfonts}
\usepackage{amssymb}
\usepackage{amsmath}
\title{1st Contest of Aether Mathematics Open}
\author{XIA Yihuan}
\date{\today}
\setlength{\parindent}{0pt}

\begin{document}

\maketitle

\section{Introduction}

Welcome to the 1st Aether Mathematics Open!
The contest will be held from September 1 to December 31, 2026.
It consists of 5 problems, each worth 20 points, for a total of 100 points. 

\paragraph{Conventions.}
Throughout this paper,
$\mathbb{N}=\{0,1,2,3,\ldots\}$ --- in accordance with the convention used in
Chinese textbooks, the natural numbers \emph{include} $0$, so $0\in\mathbb{N}$ ---
and $\mathbb{Z}^{*}=\mathbb{Z}\setminus\{0\}$;
for a subfield $K\subseteq\mathbb{C}$ and a set $S\subseteq\mathbb{C}$,
$K(S)$ denotes the smallest subfield of $\mathbb{C}$ containing $K$ and $S$;
and a polynomial is called irreducible over $K$ only if it is non-constant and
is not a product of two non-constant polynomials of $K[x]$.

\section{Problems}
\subsection{Problem 1:}

Prove that for all $n,m\in\mathbb{N}\setminus\{0,1\}$ (that is, $n,m\ge 2$)
there exist $x,y,z\in\mathbb{N}$ such that

$$x^n+y^n \equiv z^n \pmod{m}
\qquad\text{and}\qquad
z^n \not\equiv 0 \pmod{m}$$

(Here $x,y,z$ are required to be pairwise distinct.)

20 points.

\subsection{Problem 2:}

Let 

$$f(x)=\sum_{i=0}^n a_i x^i(\sum_{j=1}^n a_j^2\neq 0)\in \mathbb{Z}[x]$$

Prove that for every constant $C\in \mathbb {R}$

$${\lvert\{u \in \mathbb{R}\mid f(u) = C\}\rvert\leq n}$$

20 points.

\subsection{Problem 3:}

Let

$$K_0=\mathbb{Q}$$

$$G_i=\{p\in \mathbb{C}\mid p^o=q^r,q\in K_i,r,o,\in \mathbb{Z}^*\}$$

$$K_j=K_{j-1}(G_{j-1})$$

$$f(x)\in \mathbb{Z}[x]$$

and write $n=\deg f$ (throughout, $f$ is assumed to be non-constant).

\subsubsection{(1)}

Is it true that ($f(x)$ is irreducible over $\mathbb{R}$) $\implies$ ($f(x)$ irreducible over $K_1$)?Answer the question and prove your answer.

2 point.

\subsubsection{(2)} 

Is it true that ($f(x)$ is irreducible over $K_1$) $\implies$ ($f(x)$ is irreducible over $\mathbb{R}$)?Answer the question and prove your answer.

6 points.

\subsubsection{(3)} 

Is it true that $(\exists M \in \mathbb{N})(\forall f(x) \in \mathbb{Z}[x])[(f(x)$ irreducible over $K_M) \implies (f(x)$ irreducible over $\mathbb{R})]$?
(Recall $0\in\mathbb{N}$, so $M=0$ is allowed and then $K_M=K_0=\mathbb{Q}$.)

Determine, in terms of the degree $n=\deg f$, for which $n$ such an
implication can hold and for which it fails, and prove your answer. 

6 points.

\subsubsection{(4)} 

Let 

$$K^*_0=\mathbb{Q}$$

$$G^*_i=\{p\in \mathbb{R}\mid p^o=q^r,q\in K^*_i,r,o\in \mathbb{Z}^*\}$$

$$K^*_j=K^*_{j-1}(G^*_{j-1})$$

Find a $f(x) \in \mathbb{Z}[x]$that is irreducible over $K^*_M,(\forall M \in \mathbb{N})$but is not irreducible over $\mathbb{R}$, factor it.

6 points.

\subsubsection*{Notes}

(i) Every polynomial equation of degree $\le4$ is solvable by
radicals (the quadratic formula, Cardano, and Ferrari).

(ii) (Abel--Ruffini) There is no general formula in radicals for 
polynomial equations of degree $\ge5$; for example, no root
of $x^5-4x+2$ can be expressed by radicals. More generally,
for every $n\ge5$ the polynomial $x^n-x-1$ is irreducible over
$\mathbb{Q}$ and remains so over every field obtained from
$\mathbb{Q}$ by iteratively adjoining radicals; in particular,
none of its roots can be expressed by radicals.

(iii) (casus irreducibilis) If an irreducible cubic over
$\mathbb{Q}$ has three real roots, then none of its roots can
be expressed by real radicals, however, deeply nested.

(iv) If $g\in K_j[x]$ has degree $\le 4$, then by (i) every root
of $g$ lies in $K_{j'}$ for some $j'\ge j$. Consequently, if
$f\in\mathbb{Z}[x]$ has a factor of degree $\le4$ over some $K_j$,
then $f$ has a root that can be expressed by radicals.

\subsection{Problem 4:}

The eight moves of an ordinary chess knight can be written as the set

$$K_2=\{\pm1\pm 2i,\pm 2 \pm i\}\in\mathbb{Z}[i]$$

Now, let a "strange A-knight" be like:

$$K_a=\{\pm1\pm ai,\pm a \pm i\}\in\mathbb{Z}[i]$$

(Here $\pm1\pm ai$ stands for all four independent choices of the two signs.)

\subsubsection{(1)}:

Prove:

$$a \equiv 1 \pmod{2}\implies\sum_{i=1}^n k_i\neq 1,
\quad\text{for every } n\ge 1 \text{ and all } k_1,\ldots,k_n\in K_a$$

6 points.

\subsubsection{(2)}

Prove:

$$a \equiv 0 \pmod{2}\implies
\bigl(\forall z\in\mathbb{Z}[i]\bigr)\ \bigl(\exists n\ge 1,\ 
k_1,\ldots,k_n\in K_a\bigr):\ \sum_{i=1}^n k_i=z$$

6 points.

\subsubsection{(3)}

Let $m\ge 5$ be an integer and let

$$B_m=\{a+bi\in\mathbb{Z}[i]\mid 0\le a,b<m\}$$

(the $m\times m$ board). Prove that for every $z\in B_m$ there exist
$n\ge 1$ and $k_1,\ldots,k_n\in K_2$ such that

$$\sum_{i=1}^n k_i=z\qquad\text{and}\qquad
\sum_{i=1}^j k_i\in B_m\ \ \text{for every } 1\le j\le n ;$$

that is, an ordinary knight starting at the corner $0$ can reach every
square of the board by a path that never leaves the board.

8 points.

\subsection{Problem 5:}

Let $n\ge 1$ and let $K=(a_1,a_2,\ldots ,a_n)$ be a finite sequence with
$a_i\in\{0,1\}$. From $K$ we produce a new sequence
$T(K)=(b_1,\ldots ,b_n)$ of the same length as follows.

\paragraph{Step 1: instructions.}
Each position $i$ issues \emph{at most one} instruction.
\begin{itemize}
\item If $a_i=1$, let $p_i=\#\{j<i\mid a_j=1\}$ be the number of $1$'s
standing before it. If $p_i\ge 2$, position $i$ issues the instruction
``put a $1$ into position $i+p_i-1$''.
\item If $a_i=0$, let $s_i=\#\{j>i\mid a_j=0\}$ be the number of $0$'s
standing after it. If $s_i\ge 2$, position $i$ issues the instruction
``put a $0$ into position $i-s_i+1$''.
\end{itemize}
In every other case no instruction is issued: if the number just counted is
$0$ or $1$ (counting begins at the position itself, so the first place from
$a_i$ is $a_i$ itself), or if the indicated position does not lie in
$\{1,\ldots ,n\}$ --- there is no wrap-around and no overflow.

\paragraph{Step 2: simultaneous execution.}
All instructions are computed from the original sequence $K$ and are carried
out at the same time. If position $t$ receives only ``$1$'' instructions, put
$b_t=1$; if it receives only ``$0$'' instructions, put $b_t=0$; if it
receives both kinds, or no instruction at all, put $b_t=a_t$.

\paragraph{Question.}
Prove that $T(K)$ always has a \emph{fixed point}, i.e. there exists an index
$t\in\{1,\ldots ,n\}$ with $b_t=a_t$.

Equivalently: $T(K)$ is never the sequence $(1-a_1,1-a_2,\ldots ,1-a_n)$
obtained from $K$ by flipping every term.

20 points.

\section{Acknowledgements and Contest Notes}

The problems were proposed by XIA Yihuan, and test-solved by XIA Yihuan and some agents. 
We are grateful to Loveinterpretation for her indispensable contributions.
Teamwork is permitted; however, a team may only submit one entry. Do not use computers, calculators, large language models, or books once you fully understand the problem statements. (Large language models and machines may submit entries as well, but they are not eligible for awards; their solutions may be appended to the reference solutions if they are better.)

Everything that is necessary to solve these problems is written in the notes attached to the problems. We believe that all of you have at least completed middle school mathematics.

Please submit your solutions as a PDF to aether\_math@163.com. Both handwritten and typeset mathematical formulas are accepted; solutions presented solely in plain text without proper notation will not be accepted. Ensure that all formulas are clearly legible and mathematically correct.

Appeals and suggestions may be sent to the same email address.

\vspace{1em}
\noindent \textbf{Good luck!}

\end{document}
